\section{Introduction}\label{sec:introduction}

For \(a\in\R^3\), \(\omega\in S^2\), and \(0<\delta<1\), let
\[
T_\delta(a,\omega)=
\left\{a+t\omega+u:
 |t|\le\tfrac12,\quad u\cdot\omega=0,\quad |u|\le\delta\right\}.
\]
Thus \(T_\delta(a,\omega)\) has volume \(\pi\delta^2\). Define
\[
K_\delta f(\omega)=
\sup_{a\in\R^3}\frac1{\pi\delta^2}
\int_{T_\delta(a,\omega)}|f(x)|\,dx .
\]
The measure on \(S^2\) below is its usual surface measure.

\begin{theorem}\label{thm:main}
For every \(\varepsilon>0\) there exists a finite constant
\(C_\varepsilon\) such that, for every \(0<\delta<1\) and every
\(f\in L^3(\R^3)\),
\[
\|K_\delta f\|_{L^3(S^2)}
\le C_\varepsilon\delta^{-\varepsilon}
     \|f\|_{L^3(\R^3)}.
\]
\end{theorem}

Theorem~\ref{thm:main} resolves the three-dimensional Kakeya maximal
conjecture affirmatively. Its essential quantitative feature is uniformity
in the density of the part of each tube on which a function is large.
We explain that feature before describing the proof.

\subsection{The set problem and the maximal problem}

A Kakeya set contains a unit line segment in every direction. The problem
has its origins in the study of moving a needle inside a small planar
region; the early convex formulation appears in
\citet[Section~10]{FujiwaraKakeya1917}. Besicovitch's constructions
showed that a set with segments in every direction can have measure zero,
and that continuous needle reversal is possible in regions of arbitrarily
small area \citep{Besicovitch1928,Davies1971}. These are different
requirements. The dimension conjecture asks whether directional
containment nevertheless forces full Hausdorff dimension. Davies proved
the planar case \citep{Davies1971}; C\'ordoba's planar maximal estimates
provided a quantitative analytic counterpart~\citep{Cordoba1977}.

The distinction relevant here is already visible for characteristic
functions. Let \(\mathcal T\) be a family of unit \(\delta\)-tubes in a fixed bounded region, with
\(\delta\)-separated directions, and let a \emph{shading} assign a
measurable subset \(Y(T)\subset T\) to each tube. If
\(|Y(T)|\ge\lambda|T|\), the maximal problem asks for the cubic density
power in the union estimate
\begin{equation}\label{eq:intro-density}
\left|\bigcup_{T\in\mathcal T}Y(T)\right|
\gtrsim_\varepsilon
\delta^\varepsilon\lambda^3\sum_{T\in\mathcal T}|T|,
\qquad 0<\lambda\le1.
\end{equation}
The power loss in \(\delta\) may be arbitrarily small, uniformly in
\(\lambda\). A set-dimension conclusion can follow from weaker dependence
on \(\lambda\). In particular, Wang and Zahl's set theorem gives this
form of estimate with a power \(\lambda^{K(\varepsilon)}\); the need to
replace it by \(\lambda^3\) is stated explicitly after their theorem
\citep[Theorem~1.2]{WZ2025}. The final reduction in
Section~\ref{sec:closure} retains the density power through the passage
from shadings to arbitrary \(L^3\) functions.

The same overlapping tube configurations also enter Fourier
restriction and summability problems. Fefferman's ball-multiplier counterexample
made this connection concrete \citep{Fefferman1971}. Bourgain developed
higher-dimensional maximal estimates and their Fourier-analytic
applications \citep{Bourgain1991}, and Wolff's hairbrush argument gave maximal estimates implying
the Hausdorff-dimension bound \((n+2)/2\) in \(\R^n\)
\citep{Wolff1995}. In three dimensions this is \(5/2\). Katz,
{\L}aba, and Tao obtained a strict improvement for upper Minkowski dimension
and analyzed three recurring structures in nearly extremal tube families:
nearby tubes remain grouped across scales (stickiness), tubes meeting
near a point have nearly coplanar directions (planiness), and their
union organizes into thin planar pieces (graininess)
\citep{KLT2000}. Katz and Zahl subsequently obtained a strict
Hausdorff-dimension improvement over \(5/2\) \citep{KatzZahl2019}.

The three-dimensional set problem was resolved through a sequence of
works of Wang and Zahl. Their sticky Kakeya theorem was followed by the
full Assouad-dimension theorem and then the Hausdorff- and
Minkowski-dimension theorem for arbitrary Kakeya sets
\citep{WZSticky2026,WZAssouad2025,WZ2025}. The corresponding preprints
appeared in 2022, 2024, and 2025. The streamlined proof of Guth, Wang,
and Zahl supplies the set estimate used in this paper
\citep[Theorem~1.1]{GWZ2026}. It applies under a convex-clustering bound and a shading-density
threshold \(\lambda\ge\delta^\eta\), with \(\eta>0\) chosen from the
desired volume exponent. Convex clustering bounds
how many complete tubes can lie in a convex body; it is a formulation
that remains useful after the restrictions and changes of scale needed
here. The task is to retain the sharper density dependence required by
\eqref{eq:intro-density} through those operations. For a broader account
of the maximal formulation and its relation to the set problem, see
\citet[Conjecture~$1.3^{\prime\prime}$ and Section~3.3]{Zahl2025}.

Two further geometric estimates enter the proof. The multilinear
Kakeya theorem of Bennett, Carbery, and Tao
\citep[Theorem~1.15]{BCT2006}, sharpened to the endpoint by Guth
\citep[Theorem~1.3]{Guth2010}, controls three transverse tube families.
The planar Furstenberg theorem of Ren and Wang
\citep[Theorem~4.1]{RW2025}, in the equivalent shaded-tube formulation
recorded by Wang and Wu \citep[Theorem~0.6]{WW2024}, controls planar
incidence and projection configurations. We state these inputs and prove
the weighted forms used here in Section~\ref{sec:inputs}.

\subsection{Geometry of the proof}

\paragraph{A density-sensitive extremal problem.}
Work in a bounded coordinate chart in which each tube follows an affine
graph \(M_i(t)=b_i+t u_i\), with \(b_i,u_i\in\R^2\).
At resolution \(N=\delta^{-1}\), its shading is a set of marked time
bins. We keep the original index \(i\) and its weight under restriction,
even when two indices have the same affine graph. A full-time plank
counts the indices whose complete traces it contains; marked bins
instead determine the occupied cells and their multiplicity. This
separation permits the set input to be used after rescaling without
identifying different shadings on the same trace.

Total shading density alone does not describe what happens after a scale
cut. A temporal deficit profile records, in logarithmic units, how the
number of marked bins falls short of full branching inside each occupied
block. We introduce a family of
multiplicity inequalities whose temporal cost depends on this profile
and whose clustering term depends on the two plank widths. At the
parameter that gives the maximal theorem, the cost reduces to the total
density penalty. The critical exponent is the least extra power of
\(N\) needed to make the inequality hold for every configuration.
A positive critical exponent would produce configurations attaining
equality in the limit. We study their geometry, allowing restrictions
that lose only a subpower fraction of the weighted incidence.

\paragraph{From equality to compatible plates.}
An equality configuration cannot contain a local packet with excessive
mass: rescaling such a packet would improve the critical exponent.
The resulting local bounds, together with multilinear Kakeya, force
the directions on typical short blocks to lie close to a common plane.
The planar Furstenberg estimate supplies density inside plates containing
the reference segments. The plate Nikodym estimate then controls the
support cost of filling their missing times. Equality
therefore forces full temporal branching on an initial interval of
logarithmic scales. On those scales, the weighted set theorem captures
saturated plates: their masses attain the local bound. An anisotropic
rescaling of aligned plates rules out equality with arbitrarily small
total time deficit.

These local planes and plates are the forms of planiness and graininess
needed by the present extremal argument. Their compatibility across
scales must still be proved. We vary the two parameters in the critical
inequality to select saturated shapes, then compare two ways of cutting
an equality configuration. An outer cut removes initial scales; an
inner cut retains a fine-scale model after aligning its plates. A
minimum-delay argument forces a canonical temporal profile
\[
F(s)=\beta(s-\tau)_+,
\]
with a flat interval followed by constant deficit slope. Repeating the
outer cut reproduces this profile and the same normalized packet
geometry. This is the stationarity used below.

\paragraph{Four coordinates and two frames.}
A packet has a time interval and a nearly constant normal parameter
\(\vartheta\). In that frame, write
\[
X=M_i(t)_1+\vartheta M_i(t)_2,\qquad Y=M_i(t)_2,
\qquad U=u_{i,1}+\vartheta u_{i,2},\qquad V=u_{i,2}.
\]
Thus \(X,Y\) describe position and \(U,V\) describe velocity; the first
coordinate in each pair measures the normal component. Their admissible
widths change at four prescribed rates, determined by the stationary
time and angle scales. The masses of the packets prescribe how much
information these four coordinates must carry. Conditional entropy,
normalized by \(\log N\), measures that information as a logarithmic
cell count.

Now choose two marked events on the same original index. They give two
frames for the same affine graph. If their time and normal differences
are \(\Delta t\) and \(\Delta\vartheta\), then
\[
\begin{pmatrix}X'&Y'\\ U'&V'\end{pmatrix}
=
\begin{pmatrix}1&\Delta t\\0&1\end{pmatrix}
\begin{pmatrix}X&Y\\U&V\end{pmatrix}
\begin{pmatrix}1&0\\\Delta\vartheta&1\end{pmatrix}.
\]
Changing time adds velocity to position; changing the normal mixes the
two spatial components. The resulting scalar projections constrain the
coordinate entropy rates through the planar Furstenberg theorem.
When the two middle speeds coincide, the scale increments directly
supply only their joint rate. Freezing the middle coordinates leaves the
coefficient \(\Delta t\,\Delta\vartheta\) multiplying \(V\) in \(X'\);
this is the product slope for which the pinning estimate is needed.

The pinning argument adapts the radial-projection bootstrap of
Shmerkin and Wang \citep[Section~5.4]{SW2025} and of
Orponen, Shmerkin, and Wang \citep[Section~1.3]{OSW2024}.
On small cells, planar projection estimates improve the nonconcentration
of rays joining parameter points; this improvement can then be iterated.
Section~\ref{sec:projection-proofs} proves the finite-scale product-slope
estimate used here, including control of both single-event marginals.

The two events share an index, so their matrix and slope are generally
dependent. To apply a projection theorem, we compare their actual joint
law with the product of its conditional matrix and label marginals.
The actual law always lands in the second event's recorded support. If
a coordinate-rate constraint failed, the product law would land there
with polynomially small probability. Such a discrepancy requires a
positive amount of mutual information. An averaged short-gap estimate
shows that the available information is smaller. The resulting
projection constraints contradict the entropy demand of the stationary
packets, excluding a positive critical exponent.

\subsection{Consequences and organization}

The theorem also gives Nikodym maximal estimates in \(\R^3\) and locally
on three-dimensional manifolds of constant sectional curvature, as well
as local curved Kakeya estimates for a fixed nondegenerate translation-invariant phase
satisfying Bourgain's condition. These follow from the transfers of
Gao, Liu, and Xi \citep{GLX2025}; their full hypotheses and interpolation
ranges appear in Section~\ref{sec:consequences}. An independent route to
the same strong Kakeya maximal estimate through spherical Fourier
extension appears in \citet[Corollary~1.3]{OpenAISphere2026}.

Section~\ref{sec:critical} defines the indexed model, explains the density
reduction, and constructs the critical inequalities and equality
sequences. Section~\ref{sec:inputs} supplies the geometric inputs.
Sections~\ref{sec:local} and~\ref{sec:stationary} establish the local
geometry, canonical profile, and stationary packets.
Section~\ref{sec:stationary-estimates} proves angular nonconcentration
and the conditional normal-position estimate.

Section~\ref{sec:entropy} constructs the coordinate entropy rates,
proves their lower demand, and obtains the velocity trace needed at the
central offset. Section~\ref{sec:projections} derives the projection
constraints from the two-frame experiment;
Section~\ref{sec:projection-proofs} proves its auxiliary projection and
pinning estimates. Section~\ref{sec:closure} closes the contradiction and
passes from the density estimate to the full maximal theorem.

All limiting exponents arise from finite configurations. At each finite
stage, the required scales, partitions, and conditional tests are imposed
together before the resolution is increased. The regularization results
in Section~\ref{sec:critical} give the losses and order of limits used in
the later constructions.
